\section{Conversion of electric flux into an energy decrease}
\label{sec:flux-conversion}

This section completes the numerical argument.  The filled scalar construction
provides a metric which dominates the original metric and has the required
charged scalar-curvature inequality on a region separated from the original
boundary.  A second scalar equation converts the signed total electric flux
into a controlled decrease of ADM energy.  A lower bound for its conformal
factor retains a definite fraction of every full-cut area.  We then construct
a smooth future-trapped boundary and apply only the neutral numerical input,
Theorem~\ref{thm:neutral-input}.

\subsection{Data supplied by the filled construction}
\label{flux:setup}

Throughout this section, $n\geq4$, $k=n-1$, $\omega=\Area(S^k)$, and
$c_n^2=k(k-1)/2$.  Initially the exterior data have been prepared as in
Proposition~\ref{prep:rest-preparation} and
Lemma~\ref{prep:strictification}.  We write $g$ for their metric,
$\mathcal E_*$ for their ADM energy, and
\begin{equation}
 a_* = \inf_{\Gamma}\Area_g(\Gamma),
 \qquad X_* = (a_*/\omega)^{(k-1)/k}.
 \label{flux:prepared-area}
\end{equation}
The infimum is over the entire full cuts of this prepared exterior.  In
particular, $a_*>0$ by Lemma~\ref{prep:positive-area}.  The prepared tensor
$K$ has compact support, the metric has differentiated decay of order
$2-n$ to every fixed order, and the charged dominant-energy and prescribed
componentwise trapping inequalities are strict.  The optional compact
magnetic form is allowed in this paragraph; it will be specified again in
Corollary~\ref{flux:magnetic-extension}.

Fix $\varepsilon>0$.  For every sufficiently large approximation integer
$N$, let $\mathcal M_N$ be the filled one-ended manifold of
Proposition~\ref{pde:filled-solve}.  We also denote its background metric by
$g$.  It agrees with the prepared metric on the original exterior.  Put
\begin{equation}
 \ell=e^{-N\varepsilon},\qquad \eta=\ell^{3/2},\qquad
 h=\eta f,\qquad z=t+\varepsilon,\qquad
 \bar g=g+l^2\,\dd f^2,\qquad g_b=e^{2z}\bar g.
 \label{flux:background-metrics}
\end{equation}
The scalar-construction variables satisfy
\begin{equation}
 z\geq0,\qquad \ell\leq l\leq e,\qquad
 |\dd h|_{\bar g}\leq\eta/l\leq\ell^{1/2}.
 \label{flux:height-gradient}
\end{equation}
For fixed prepared data and $\varepsilon$, the bounds for $t$, $Z$, and
$\nabla f$ imply that the metric coefficients and inverse coefficients
needed below have bounds of the form $\exp(\mathcal P_N)$, where
$\mathcal P_N$ denotes a polynomial in $1+N$ whose coefficients may change
from occurrence to occurrence.  No uniform bound as $N\to\infty$ is claimed
for the second scalar solution.

We use the following further conclusions of
Propositions~\ref{pde:mass-cost}, \ref{pde:height-separation}, and
\ref{pde:charged-region}.  There are constants $b_2>b_3>0$, independent of
$N$, such that
\begin{equation}
 |h|>b_2\quad\hbox{on }S,\qquad h\longrightarrow0
 \quad\hbox{at infinity},
 \label{flux:height-separation}
\end{equation}
and the set $\{|h|\geq b_3/4\}$ on the original exterior is contained in a
fixed compact set.  We take $b_3<b_2/2$.  On
\begin{equation}
 \mathcal U_N=\{x\hbox{ in the original exterior}: |h(x)|\leq b_3\},
 \label{flux:low-region}
\end{equation}
the retained curvature estimate is the charged one, with an additional
positive margin on $b_3/4\leq|h|\leq b_3$.  More precisely, the scalar
identity supplies there a constant $c_*>0$, independent of $N$, and the
nonnegative term $(1-\delta_0)\mathcal T$, with $\delta_0<1/2$.
The energy estimate, in coordinates normalized for $g_b$ at infinity, is
\begin{equation}
 \mathcal E(g_b)\leq
 e^{(k-1)\varepsilon}(\mathcal E_*+o_N(1)),
 \qquad o_N(1)\longrightarrow0
 \quad(N\longrightarrow\infty).
 \label{flux:background-energy}
\end{equation}
The fixed-$N$ metric $g_b$ has all-order differentiated asymptotic decay
of order $2-n$ after the constant coordinate dilation, and
$R_{g_b}=O(r^{-n-\beta})$ for a $\beta>0$.

Let $\alpha=i_E\vol_g$ on the original exterior.  Multiply it by
$\sgn Q$ when $Q\neq0$, and use either sign when $Q=0$.  Extend this signed
form smoothly into the filling, cutting it off on fixed collars on the
inside of $S$; call the extension $\widetilde\alpha$.  It is closed on the
original exterior and need not be closed in the filling.  Define
\begin{equation}
 i_{E_b}\vol_{g_b}=\widetilde\alpha,
 \qquad i_{E_{\mathrm{ext}}}\vol_g=\widetilde\alpha.
 \label{flux:extended-fields}
\end{equation}
Thus the flux of $E_b$ across a large sphere is $\omega|Q|$.  Its divergence
vanishes on the original exterior.  The vector $E_{\mathrm{ext}}$ is
uniformly bounded on the filling and is $O(r^{-k})$ on the end.  Its support
in the filling can be kept in the fixed collars; even an extension over the
whole core would have only polynomially growing $L^p$ norms, because that
core has volume $O(1+N)$.

Matrix formulas below use a $g$-orthonormal frame, or equivalently the
$g$-identification of vectors and covectors.  For clarity, the form norm
underlying the charged estimate is
\begin{equation}
 |\alpha|_{\bar g}^2
   =(1-|w|_g^2)|E|_g^2+g(E,w)^2
   =|E|_g^2-|w\wedge E|_g^2.
 \label{flux:graph-form-norm}
\end{equation}
Indeed $\bar g^{-1}=A_d=g^{-1}-w\otimes w$ and
$\vol_{\bar g}=\sqrt D\,\vol_g$.  Completing the square in the magnetic
two-form, when it is present, gives
$I_g(w)\geq c_n^2|\alpha|_{\bar g}^2$.  Since $z\geq0$,
\begin{equation}
 e^{2z}|E_b|_{g_b}^2
 =e^{-2(k-1)z}|\alpha|_{\bar g}^2
 \leq |\alpha|_{\bar g}^2.
 \label{flux:conformal-field-cost}
\end{equation}
Consequently the scalar-construction estimates imply
$R_{g_b}/2\geq c_n^2|E_b|_{g_b}^2$ throughout $\mathcal U_N$.

\begin{lemma}[A retained gradient term]
\label{flux:retained-gradient}
The quadratic expression $\mathcal T$ in the filled scalar construction
satisfies
\begin{equation}
 \mathcal T\geq\frac{k^2-1}{4}|\dd t|_{\bar g}^2.
 \label{flux:T-gradient}
\end{equation}
In particular, on the original band $b_3/4\leq|h|\leq b_3$,
\begin{equation}
 e^{2z}R_{g_b}/2\geq
 c_n^2 e^{2z}|E_b|_{g_b}^2+c_*+
 \frac{k^2-1}{8}|\dd t|_{\bar g}^2.
 \label{flux:background-band-curvature}
\end{equation}
\end{lemma}

\begin{proof}
The inverse graph metric is $A_d$. Hence
\eqref{flux:T-gradient} is exactly
\eqref{alg:gradient-coercivity} in the Coercivity Lemma
(Lemma~\ref{alg:coercivity-lemma}). Combining that estimate with
$1-\delta_0\geq1/2$ in \eqref{pde:retained-band}, the electromagnetic
square completion above, and \eqref{flux:conformal-field-cost}
gives \eqref{flux:background-band-curvature}.
\end{proof}

\subsection{Smooth profiles with the sharp limiting coefficient}

\begin{lemma}[Flux profiles]
\label{flux:profiles}
Fix $0<s<1$ and set
\begin{equation}
 y_0=\frac{\log(1-s^2)}{k-1}<0.
 \label{flux:floor}
\end{equation}
For every $\delta>0$ there exist smooth functions $L>0$ and $H_0$ on
$\mathbb R$ with
\begin{gather}
 L(0)=1,\qquad L'\geq0,\qquad
 H_0=0\text{ on }(-\infty,y_0],\qquad
 \supp H_0'\Subset(y_0,0),\qquad H_0'\geq0,
 \label{flux:profile-properties}\\
 |H_0'(y)|\leq
 k\sqrt{(k-1)\bigl(k-1+2L'(y)/L(y)\bigr)}
 e^{(k-1)y}L(y),
 \label{flux:profile-discriminant}\\
 ks-\delta<H_0(0)\leq ks.
 \label{flux:profile-value}
\end{gather}
In particular $H_0$ is bounded.  The function defined by
\begin{equation}
 B_0(0)=0,\qquad B_0'(y)=e^{(k-1)y}L(y)
 \label{flux:B-definition}
\end{equation}
is strictly increasing, and its range on $[y_0,\infty)$ is
$[B_0(y_0),\infty)$.
\end{lemma}

\begin{proof}
Let $a_0=k-1$ and $c=1-s^2$.  On $(y_0,\infty)$ consider
\begin{equation}
 L_{\mathrm{opt}}(y)^2=
 \frac{1-c e^{-a_0y}}{s^2}.
 \label{flux:optimal-profile}
\end{equation}
This is positive and increasing, and $L_{\mathrm{opt}}(0)=1$.  Direct
differentiation gives
\begin{equation}
 k\sqrt{a_0(a_0+2L_{\mathrm{opt}}'/L_{\mathrm{opt}})}
 e^{a_0y}L_{\mathrm{opt}}=
 \frac{ka_0}{s}e^{a_0y}.
 \label{flux:optimal-derivative}
\end{equation}
Its integral from $y_0$ to zero is $ks$.

Here is a positive smooth replacement at the lower endpoint.  Take a small
$\sigma>0$ and a smooth nondecreasing cutoff $\chi_\sigma$ which is zero
below $y_0+\sigma/3$ and one above $y_0+2\sigma/3$.  For
$y\leq b_\sigma:=y_0+\sigma$ define
\begin{equation}
 L(y)=L_{\mathrm{opt}}(b_\sigma)
 -\int_y^{b_\sigma}
   \chi_\sigma(v)L_{\mathrm{opt}}'(v)\,\dd v,
 \label{flux:profile-smoothing}
\end{equation}
where the integrand is zero below $y_0+\sigma/3$.  For $y\geq b_\sigma$
use $L=L_{\mathrm{opt}}$.  The definitions agree with all derivatives near
$b_\sigma$.  The lower constant value is positive, since it is at least
$L_{\mathrm{opt}}(y_0+\sigma/3)$.  Thus $L$ is smooth, positive and
nondecreasing on the whole line, and still has $L(0)=1$.

Choose $0\leq\zeta_\sigma\leq1$ smooth with compact support in
$(y_0+\sigma,-\sigma)$, tending pointwise to one on $(y_0,0)$ as
$\sigma\downarrow0$, and put
\begin{equation}
 H_0'(y)=\zeta_\sigma(y)\frac{ka_0}{s}e^{a_0y},
 \qquad H_0(y)=\int_{-\infty}^y H_0'(v)\,\dd v.
 \label{flux:H-construction}
\end{equation}
On the support of $H_0'$ the function $L$ is exactly
$L_{\mathrm{opt}}$, so \eqref{flux:profile-discriminant} follows from
\eqref{flux:optimal-derivative}.  Elsewhere its left side is zero.
Dominated convergence gives $H_0(0)\to ks$ from below, proving
\eqref{flux:profile-value} after decreasing $\sigma$.
Finally, $B_0'>0$ and, for $y\geq0$, $L(y)\geq1$, so
\begin{equation}
 B_0(y)\geq\frac{e^{(k-1)y}-1}{k-1}.
 \label{flux:B-growth}
\end{equation}
This proves the range assertion.
\end{proof}

\subsection{The bounded-flux equation on the filled manifold}

\begin{proposition}[Filled flux conversion]
\label{flux:conversion-solve}
Fix the prepared data, $\varepsilon$, $s$, and one smooth profile from
Lemma~\ref{flux:profiles}.  For each sufficiently large fixed $N$ there is
a smooth function $y$ on $\mathcal M_N$ such that
\begin{equation}
 k\Delta_{g_b}B_0(y)=\operatorname{div}_{g_b}(H_0(y)E_b),
 \qquad y\longrightarrow0\quad\hbox{at infinity}.
 \label{flux:conversion-equation}
\end{equation}
It satisfies
\begin{equation}
 y\geq y_0,\qquad y\leq\mathcal P_N,
 \qquad y=O_j(r^{2-n})\quad\hbox{for every fixed }j,
 \label{flux:y-bounds}
\end{equation}
where constants in the last estimates may depend on $N$ and on the fixed
profile.  Moreover $B_0(y)$ is harmonic sufficiently far out, and
\begin{equation}
 \lim_{R\to\infty}\int_{S_R}\partial_{\nu_b}y\,\dd A_{g_b}
 =\frac{H_0(0)}{k}\,\omega|Q|.
 \label{flux:y-flux}
\end{equation}
\end{proposition}

\begin{proof}
We give the estimates on expanding smooth coordinate-sphere truncations of
$\mathcal M_N$, with zero outer Dirichlet value.  There is no inner boundary.

\paragraph{A background Sobolev inequality.}
The background filled metric has uniformly controlled local geometry.  Its
core is covered by at most $C(1+N)$ coordinate patches of fixed size and
bounded overlap; the product necks account for the number of patches.
The overlap graph is connected to a fixed annulus in the unchanged
asymptotic end by chains of length at most $C(1+N)$.  For a compactly
supported smooth function $v$, radial integration from infinity and
Cauchy--Schwarz give, on this fixed annulus $A$,
\begin{equation}
 \|v\|_{L^2(A,g)}^2\leq
 C\int_{\mathrm{end}}|\dd v|_g^2\,\vol_g.
 \label{flux:anchor-annulus}
\end{equation}
For example, the radial weight in this estimate is finite because
$\int_R^\infty r^{1-n}\,\dd r<\infty$.  Local Poincar\'e inequalities
control the oscillation about each patch mean.  On overlapping patches the
difference of the two means is controlled by the gradient energy on their
union.  Summing along the chains to $A$, then summing over the patches,
therefore gives
\begin{equation}
 \|v\|_{L^2(\mathrm{core},g)}^2
 \leq C(1+N)^C\int_{\mathcal M_N}|\dd v|_g^2\,\vol_g.
 \label{flux:core-poincare}
\end{equation}
Combine this with the local Sobolev inequalities on the patches and the
ordinary Euclidean Sobolev inequality on the end, using a cutoff across
$A$.  With $p_*=2n/(n-2)$, this proves
\begin{equation}
 \|v\|_{L^{p_*}(g)}^2\leq
 S_N\int_{\mathcal M_N}|\dd v|_g^2\,\vol_g,
 \qquad S_N\leq\exp(\mathcal P_N).
 \label{flux:background-Sobolev}
\end{equation}
The polynomial bound implicit above would suffice.  The weaker displayed
bound is convenient for combining constants.  Functions with zero trace
on a truncation can be extended by zero, so the same constant works for
all outer radii.  The local Sobolev and Poincar\'e inequalities used here are
the usual uniformly elliptic coordinate versions; see
\cite[Chapters 7--8]{GilbargTrudinger2001}.

\paragraph{The exact divergence matrix.}
Set $Y=B_0(y)$ and $Y_0=B_0(y_0)<0$.  Define
\begin{equation}
 \mathfrak h(Y)=H_0(B_0^{-1}(Y))\quad(Y\geq Y_0),\qquad
 \mathfrak h(Y)=0\quad(Y<Y_0).
 \label{flux:bounded-coefficient}
\end{equation}
This is smooth and bounded on the whole real line.  It is identically zero
near and below $Y_0$, and constant for all sufficiently large $Y$; all its
nonzero derivatives have compact support.  The ratio of volume forms and
the inverse metric are
\begin{equation}
 J_b=\frac{\vol_{g_b}}{\vol_g}=e^{nz}\sqrt D,
 \qquad g_b^{-1}=e^{-2z}A_d.
 \label{flux:volume-and-inverse}
\end{equation}
Consequently the equation in the fixed background measure is exactly
\begin{equation}
 k\operatorname{div}_g(A\nabla_gY)
   =\operatorname{div}_g(\mathfrak h(Y)E_{\mathrm{ext}}),
 \qquad A=e^{(n-2)z}\sqrt D\,A_d.
 \label{flux:background-equation}
\end{equation}
Here $J_bE_b=E_{\mathrm{ext}}$ follows directly from
\eqref{flux:extended-fields}.  In particular no graph-volume factor remains
in the right-hand flux.

The transverse eigenvalues of $A$ are $e^{(n-2)z}\sqrt D$ and its axial
eigenvalue is $e^{(n-2)z}/\sqrt D$.  The filled-solve estimates therefore give
\begin{equation}
 \lambda_N\,\id\leq A\leq\Lambda_N\,\id,
 \qquad \lambda_N^{-1}+\Lambda_N\leq\exp(\mathcal P_N).
 \label{flux:ellipticity}
\end{equation}
If $F_Y=\mathfrak h(Y)E_{\mathrm{ext}}$, its $L^2(g)$ and $L^{2n}(g)$ norms have
polynomial bounds in $1+N$, uniformly in $Y$ and in the truncation radius.
Indeed $\mathfrak h$ is bounded, the core extension was just described, and the end
integral of $r^{-pk}$ is finite for $p=2,2n$.

\paragraph{Energy and the level estimate.}
The same estimates apply to the homotopy which multiplies the right-hand
side of \eqref{flux:background-equation} by a number in $[0,1]$.
Testing by $Y$ gives
\begin{equation}
 k\lambda_N\|\dd Y\|_2^2
 \leq\|F_Y\|_2\|\dd Y\|_2,
 \qquad \|Y\|_{p_*}\leq\exp(\mathcal P_N).
 \label{flux:energy-bound}
\end{equation}
For $s_1>0$, put $u_{s_1}=(Y-s_1)_+$ and
$A_{s_1}=\{Y>s_1\}$.  Testing by $u_{s_1}$, applying H\"older with exponent
$2n$, and then using \eqref{flux:background-Sobolev}, yields
\begin{equation}
 \|u_{s_1}\|_{p_*}
 \leq C_N |A_{s_1}|^{1/2-1/(2n)},
 \qquad C_N\leq\exp(\mathcal P_N).
 \label{flux:level-norm}
\end{equation}
Thus, for $t_1>s_1$,
\begin{equation}
 |A_{t_1}|\leq
 \left(\frac{C_N}{t_1-s_1}\right)^{p_*}
 |A_{s_1}|^{\vartheta},
 \qquad \vartheta=\frac{n-1}{n-2}>1.
 \label{flux:level-recursion}
\end{equation}
This also holds for the negative levels of $Y$.

For completeness, start at level one, whose superlevel measure is at most
$\exp(\mathcal P_N)$ by \eqref{flux:energy-bound}, and set
$s_j=1+d_*(1-2^{-j})$.  The recursion has the form
$M_{j+1}\leq C_N^{p_*}d_*^{-p_*}2^{p_*(j+1)}M_j^\vartheta$.
Choosing $d_*$ to be a fixed sufficiently large multiple of
$C_N M_0^{(\vartheta-1)/p_*}$ makes $M_j$ tend to zero geometrically after
a further fixed exponential rescaling in $j$.  All the constants and the
starting measure are bounded by $\exp(\mathcal P_N)$, so this choice has
the same bound.  Therefore
\begin{equation}
 \|Y\|_\infty\leq\exp(\mathcal P_N)
 \label{flux:Y-supremum}
\end{equation}
uniformly in the radius and in the homotopy parameter.

\paragraph{Existence on a truncation and exhaustion.}
Fix $N$ and an outer radius.  For a continuous trial function $Y$ of zero
boundary value, freeze $\mathfrak h(Y)E_{\mathrm{ext}}$ and solve the linear
Dirichlet equation \eqref{flux:background-equation}.  On this fixed smooth
domain its coefficients are smooth and uniformly elliptic.  Linear
divergence-form $W^{1,p}$ estimates with any fixed $p>n$ show that this
solution operator is continuous and compact into continuous functions of
zero boundary value.  Fixed points of its homotopy with the zero map
satisfy \eqref{flux:Y-supremum}.  Leray--Schauder degree on a larger ball
therefore gives a fixed point, since the degree at the zero map is one;
see \cite[Sections 12--17]{LeraySchauder1934}.

For such a solution the bounded divergence data first give local H\"older
estimates.  Then $\mathfrak h(Y)E_{\mathrm{ext}}$ is H\"older continuous, divergence
Schauder estimates give $C^{1,\alpha}$ regularity, and successive
differentiation gives smoothness.  These estimates, with the supremum
bound, are uniform on each fixed compact subset as the outer radius tends
to infinity.  The standard interior and Dirichlet estimates used at this
step apply to smooth uniformly elliptic scalar operators on fixed charts;
see \cite[Chapters 6 and 8]{GilbargTrudinger2001}.  A diagonal subsequence
thus gives a smooth solution on $\mathcal M_N$, preserving the energy and
supremum bounds.

\paragraph{The floor and inversion.}
On a truncation, test \eqref{flux:background-equation} by
$(Y_0-Y)_+$.  Its support is a region on which $\mathfrak h(Y)=0$.  Ellipticity gives
zero gradient for this test function, and its outer boundary trace is
zero because $Y_0<0$.  Hence $Y\geq Y_0$.  This passes to the limit.
Lemma~\ref{flux:profiles} now defines $y=B_0^{-1}(Y)$ everywhere.  The
inverse is smooth at $Y_0$ as well, since $B_0'(y_0)>0$.  Inequality
\eqref{flux:B-growth} and \eqref{flux:Y-supremum} give
$y\leq\mathcal P_N$, while the floor is $y\geq y_0$.

\paragraph{The end and the total flux.}
On the end $\operatorname{div}_{g_b}E_b=0$, so the equation for $Y$ is
\begin{equation}
 k\Delta_{g_b}Y=\mathfrak h'(Y)E_b(Y).
 \label{flux:exterior-drift}
\end{equation}
The drift coefficient is $O(r^{-k})$, with fixed-$N$ bounds.  Choose
$0<\beta'<\min(1,n-2)$.  After increasing the fixed inner radius, positive
multiples of $r^{2-n}(1-r^{-\beta'})$ are supersolutions of the absolute
drift comparison problem: their leading negative Laplacian term has
order $r^{-n-\beta'}$, and dominates the metric-error and first-order
terms.  Comparison with both signs, using the uniform bound on the fixed
inner sphere and zero data at the truncation sphere, gives
$Y=O(r^{2-n})$ uniformly during exhaustion.  Scaled elliptic estimates
for \eqref{flux:exterior-drift} then give the first two differentiated
orders, and the same is true for $y$.

Since $\supp H_0'\Subset(y_0,0)$ and $y\to0$, one has
$H_0(y)=H_0(0)$ identically sufficiently far out.  There $Y$ is harmonic
for $g_b$.  The all-order metric estimates for $g_b$ and scaled harmonic
estimates yield $Y,y=O_j(r^{2-n})$ for every fixed $j$.
Integrating \eqref{flux:conversion-equation} on the whole filled
truncation now gives, for every sufficiently large radius,
\begin{equation}
 k\int_{S_R}\partial_{\nu_b}Y\,\dd A_{g_b}
 =H_0(0)\int_{S_R}g_b(E_b,\nu_b)\,\dd A_{g_b}
 =H_0(0)\omega|Q|.
 \label{flux:Y-flux}
\end{equation}
There is no omitted inner flux.  Because $B_0'(0)=1$ and
$B_0'(y)-1=O(y)$, replacing $\dd Y$ by $\dd y$ in this integral has an
error $O(R^{2-n})$.  This proves \eqref{flux:y-flux}.
\end{proof}

\subsection{Curvature, mass, and a smooth trapped boundary}

\begin{lemma}[The uncharged metric]
\label{flux:uncharged-metric}
Let $y$ be supplied by Proposition~\ref{flux:conversion-solve} and set
\begin{equation}
 \gamma=e^{2y}g_b=e^{2(z+y)}\bar g.
 \label{flux:gamma}
\end{equation}
Then $R_\gamma\geq0$ on $\mathcal U_N$ and
\begin{equation}
 \gamma\geq e^{2y_0}g.
 \label{flux:metric-dominance}
\end{equation}
Its ADM energy is
\begin{equation}
 \mathcal E(\gamma)=\mathcal E(g_b)-\frac{H_0(0)}{k}|Q|.
 \label{flux:mass-decrease}
\end{equation}
On the band $b_3/4\leq|h|\leq b_3$ in the original exterior,
\begin{equation}
 e^{2(z+y)}R_\gamma/2\geq
 c_*+\frac{k^2-1}{8}|\dd t|_{\bar g}^2+
 \mathbf 1_{\{y\geq0\}}\frac{k(k-1)}2|\dd y|_{\bar g}^2.
 \label{flux:gamma-band-curvature}
\end{equation}
\end{lemma}

\begin{proof}
Where the extended flux form is closed, expand the equation as
\begin{equation}
 k\Delta_{g_b}y=
 \frac{H_0'}{B_0'}E_b(y)
 -k\left(k-1+\frac{L'}L\right)|\dd y|_{g_b}^2.
 \label{flux:expanded-y-equation}
\end{equation}
The conformal scalar formula gives
\begin{align}
 e^{2y}R_\gamma/2
 &=R_{g_b}/2-k\Delta_{g_b}y
       -\frac{k(k-1)}2|\dd y|_{g_b}^2\notag\\
 &\geq c_n^2|E_b|_{g_b}^2
 -\frac{H_0'}{e^{(k-1)y}L}E_b(y)
 +k\left(\frac{k-1}{2}+\frac{L'}L\right)|\dd y|_{g_b}^2.
 \label{flux:curvature-quadratic}
\end{align}
The last expression is nonnegative: the square of its mixed coefficient
is at most
\begin{equation}
 4\frac{k(k-1)}2\,
 k\left(\frac{k-1}{2}+\frac{L'}L\right)
 =k^2(k-1)\left(k-1+2\frac{L'}L\right),
 \label{flux:discriminant-check}
\end{equation}
exactly by \eqref{flux:profile-discriminant}.  Completing the square in
$E_b$ proves the assertion on $\mathcal U_N$.  The positive terms beyond
the electric term in \eqref{flux:background-band-curvature} have not been
used in this completion.  If $y\geq0$, then $H_0'=0$ and $L'\geq0$, so the
entire additional term $k(k-1)|\dd y|_{g_b}^2/2$ can also be retained.
Multiplication by $e^{2z}$ proves \eqref{flux:gamma-band-curvature}.
The metric comparison follows from $z\geq0$, $y\geq y_0$, and
$\bar g\geq g$.

In coordinates normalized for $g_b$, the conformal ADM calculation is
\begin{equation}
 \mathcal E(e^{2y}g_b)-\mathcal E(g_b)
 =-\frac1\omega\lim_{R\to\infty}
       \int_{S_R}\partial_{\nu_b}y\,\dd A_{g_b}.
 \label{flux:conformal-energy-formula}
\end{equation}
To check the normalization, the leading perturbation is $2y\delta_{ij}$;
its Euclidean ADM integrand is $-2k\partial_i y$.  The normalizing factor
is $1/(2k\omega)$.  The difference between geometric and Euclidean
surface fluxes, and every quadratic error, is $O(R^{2-n})$.  Equation
\eqref{flux:y-flux} now proves \eqref{flux:mass-decrease}.
The end scalar curvature remains integrable: outside a compact set $Y$ is
harmonic, so \eqref{flux:expanded-y-equation} has no mixed term and
$\Delta_{g_b}y=O(r^{2-2n})$, while $R_{g_b}=O(r^{-n-\beta})$.
\end{proof}

\begin{lemma}[A freely chosen late plateau]
\label{flux:trace-tensor}
For all sufficiently large $N$ one can choose a smooth nonnegative
function $\mathcal H$ of $|h|$, zero on $[0,b_3/4]$ and equal to any
prescribed finite nonnegative constant on $[b_3/2,b_3]$, so that the tensor
\begin{equation}
 K_{\mathrm{new}}=-\mathfrak a\,\gamma,
 \qquad
 \mathfrak a=e^{-(z+b(y))}\mathcal H(|h|)
 \label{flux:new-tensor}
\end{equation}
satisfies the neutral dominant energy condition on $\mathcal U_N$.
Here $b$ is a fixed smooth function with
\begin{equation}
 b=0\text{ on }(-\infty,0],\qquad 0\leq b'\leq1,
 \qquad y-b(y)\text{ bounded on }[y_0,\infty).
 \label{flux:b-cutoff}
\end{equation}
The lower bound on $N$ is independent of the prescribed plateau height.
\end{lemma}

\begin{proof}
Choose $b'$ smooth, zero up to zero and equal to one eventually, and
integrate from zero.  Then
\begin{equation}
 0<C_-\leq C(y):=e^{y-b(y)}\leq C_+<\infty
 \quad(y\geq y_0),
 \label{flux:C-bounds}
\end{equation}
with constants independent of $N$.
For a pure-trace tensor $-\mathfrak a\gamma$ the neutral densities are
\begin{equation}
 \mu_{\mathrm{new}}=R_\gamma/2+\frac{nk}{2}\mathfrak a^2,
 \qquad J_{\mathrm{new}}=k\,\dd\mathfrak a.
 \label{flux:trace-constraints}
\end{equation}
On the band, multiply the required inequality by $e^{2(z+y)}$.
The new quadratic energy is $nk C^2\mathcal H^2/2$, while
\eqref{flux:height-gradient} bounds the momentum by
\begin{equation}
 kC\left(\mathcal H|\dd t+b'\dd y|_{\bar g}
              +(\eta/l)|\mathcal H'|\right).
 \label{flux:trace-momentum-bound}
\end{equation}
Set $A_t=(k^2-1)/8$ and $B_y=k(k-1)/2$.  Young's inequality absorbs the
first term of \eqref{flux:trace-momentum-bound} into the two gradient
terms in \eqref{flux:gamma-band-curvature}, at a cost at most
$C^2\mathcal H^2\chi_k$, where
\begin{equation}
 \chi_k=\frac{k^2}{4}\left(\frac1{A_t}+\frac1{B_y}\right)
 =\frac{2k^2}{k^2-1}+\frac{k}{2(k-1)}\leq3
 \quad(k\geq3).
 \label{flux:dimension-slack}
\end{equation}
There is no use of the $y$-gradient term where $y<0$, because $b'=0$
there.  Where $y\geq0$, that term was retained in full in the preceding
lemma.  Hence no curvature term is used twice.  As $nk/2\geq6$, the
remaining coefficient
\begin{equation}
 d_k=\frac{nk}{2}-\chi_k\geq3
 \label{flux:positive-slack}
\end{equation}
is positive in every dimension under consideration.

It is now enough that
\begin{equation}
 (\eta/l)|\mathcal H'|\leq c'(1+\mathcal H^2),
 \qquad
 0<c'\leq
 \min\left\{\frac{c_*}{kC_+},\frac{d_kC_-}{k}\right\}.
 \label{flux:transition-condition}
\end{equation}
Indeed this makes the second term in \eqref{flux:trace-momentum-bound}
at most $c_*+d_kC^2\mathcal H^2$.
Choose a smooth cutoff $\chi_h$ which is zero on $[0,b_3/4]$ and one on
$[b_3/2,b_3]$.  For the desired plateau $M<\infty$, define
\begin{equation}
 \mathcal H(s_1)=\tan\bigl((\arctan M)\chi_h(s_1)\bigr).
 \label{flux:arctan-ramp}
\end{equation}
Then
\begin{equation}
 \frac{|\mathcal H'|}{1+\mathcal H^2}
 \leq\frac\pi2\|\chi_h'\|_\infty,
 \label{flux:ramp-uniformity}
\end{equation}
independently of $M$.  Since $\eta/l\leq\ell^{1/2}\to0$, one sufficiently
large $N$ ensures \eqref{flux:transition-condition} for every finite
plateau.  Off the band the tensor vanishes on the required exterior and
$R_\gamma\geq0$.  Although $|h|$ need not be smooth at its zero set, its
composition with $\mathcal H$ is smooth because $\mathcal H$ vanishes on a
neighborhood of zero.  This proves the lemma.
\end{proof}

Figure~\ref{flux:fig-height-bands} records the height bands and the
order of the remaining choices.
\begin{figure}[tbp]
\centering
\begin{tikzpicture}[x=1cm,y=1cm,font=\small,
  line cap=round,line join=round,
  every node/.style={inner sep=2pt},
  axis/.style={draw=black!65,line width=.5pt},
  band/.style={draw=black!55,fill=black!8,line width=.45pt},
  focus/.style={draw=linkblue,fill=linkblue!9,line width=.6pt}]
  \node[draw=black!40,fill=black!3,anchor=north west,
    text width=5.25cm,minimum height=1.22cm,align=left]
    at (0,6.4) {\textbf{Added region}\\
      No energy condition imposed.\\
      $\widetilde\alpha$ supported in fixed inner collars.};
  \node[draw=black!40,fill=black!3,anchor=north west,
    text width=9.05cm,minimum height=1.22cm,align=left]
    at (6.25,6.4) {\textbf{Original exterior}\\
      $\widetilde\alpha$ is closed.\quad
      On $S$: $|h|>b_2>2b_3$.\\
      At the distinguished end: $h\longrightarrow0$.};
  \node[anchor=west] at (0,4.78)
    {Both scalar solves take place on the filled manifold $\mathcal M_N$.};

  \draw[axis,-{Latex[length=2mm]}] (5,3.88)--(15.45,3.88);
  \node[anchor=west] at (15.5,3.88) {$|h|$};
  \foreach \x/\lab in {5/0,7.25/{b_3/4},9.5/{b_3/2},11.4/{c_h},14/{b_3}} {
    \draw[axis] (\x,3.79)--(\x,3.97);
    \node[anchor=south] at (\x,4.02) {$\lab$};
  }
  \draw[draw=linkblue,densely dashed,line width=.7pt]
    (11.4,.43)--(11.4,3.8);
  \node[anchor=west] at (0,3.25) {Charged estimate on $\mathcal U_N$};
  \draw[band] (5,2.98) rectangle (14,3.52);
  \node at (9.5,3.25) {$R_{g_b}/2\geq c_n^2|E_b|_{g_b}^2$};

  \node[anchor=west] at (0,2.45) {Additional positive margin};
  \draw[focus] (7.25,2.18) rectangle (14,2.72);
  \node at (10.625,2.45) {$b_3/4\leq |h|\leq b_3$};

  \node[anchor=west] at (0,1.65) {Trace profile $\mathcal H(|h|)$};
  \draw[band] (5,1.38) rectangle (7.25,1.92);
  \draw[band,fill=black!14] (7.25,1.38) rectangle (9.5,1.92);
  \draw[band,fill=black!22] (9.5,1.38) rectangle (14,1.92);
  \node at (6.125,1.65) {$0$};
  \node at (8.375,1.65) {transition};
  \node at (11.75,1.65) {plateau $M$};

  \node[anchor=west] at (0,.85) {Heights in the retained $\mathcal D_N$};
  \draw[focus] (5,.58) rectangle (11.4,1.12);
  \node at (8.2,.85) {$0\leq |h|\leq c_h$};

  \node[anchor=north west,align=left,text width=15.85cm] at (0,.05)
    {$\displaystyle
      \mathcal D_N=\overline{\text{component of }\{|h|<c_h\}
      \text{ containing the far end}},\qquad b_3/2<c_h<b_3.$\\[3pt]
     $\displaystyle
      \supp(K_{\mathrm{new}}|_{\mathcal D_N})
      \subset\{x\in\mathcal D_N:b_3/4\leq |h(x)|\leq c_h\}.$};
\end{tikzpicture}
\caption{Height bands and the late neutral cut. The bars refer only to the
original exterior and encode height, not distance or the topology of level
sets. The end component $\mathcal D_N$ stays outside the filling; every
component of its boundary is retained. After the two scalar solves and the
regular choice of $c_h$, the finite plateau $M$ is chosen to make that
boundary strictly future trapped, without changing $N$ or either solution.
See \eqref{flux:height-separation}--\eqref{flux:low-region},
Lemma~\ref{flux:trace-tensor}, and \eqref{flux:strict-new-boundary}.}
\label{flux:fig-height-bands}
\end{figure}


\begin{proposition}[Prepared numerical bound]
\label{flux:prepared-bound}
For the prepared data, for every $0<s<1$,
\begin{equation}
 \mathcal E_*\geq s|Q|+\frac{1-s^2}{2}X_*.
 \label{flux:prepared-inequality}
\end{equation}
The conclusion also holds for prepared data carrying a smooth compactly
supported closed magnetic two-form, with the charged energy convention
used in the filled scalar construction.
\end{proposition}

\begin{proof}
Keep $s$, a smooth profile, and $\varepsilon$ fixed.  Take $N$ sufficiently
large for all preceding constructions.  Sard's theorem gives a regular
value $c_h\in(b_3/2,b_3)$ of $|h|$.  Let $\mathcal D_N$ be the closure of
the component of $\{|h|<c_h\}$ containing the entire far end.
By \eqref{flux:height-separation}, this component cannot cross $S$ into the
filling.  Its entire boundary is a compact smooth embedded two-sided
hypersurface in the original interior.  It is nonempty: a path in the
original connected exterior from the far end to $S$ must cross the chosen
level.  All boundary components are retained.

The mean curvature of this fixed compact boundary in $\gamma$ is finite.
We may now choose the plateau $M$ in Lemma~\ref{flux:trace-tensor} so large
that, with normal toward the end,
\begin{equation}
 H_{\partial\mathcal D_N}
       +\operatorname{tr}_{\partial\mathcal D_N}K_{\mathrm{new}}
 =H_{\partial\mathcal D_N}-k\mathfrak a<0
 \quad\hbox{on every component}.
 \label{flux:strict-new-boundary}
\end{equation}
This late choice does not require repeating either scalar solve or
increasing $N$, by \eqref{flux:ramp-uniformity}.  Restricted to
$\mathcal D_N$, the new tensor has compact support, since $h\to0$ and
$\mathcal H$ vanishes near zero.

We check the neutral input's hypotheses explicitly.  The exterior
$\mathcal D_N$ is connected and oriented, smooth up to its nonempty compact
boundary, and has exactly the original coordinate end with compact
complement.  It is complete with boundary included: it is closed in the
original complete exterior, and \eqref{flux:metric-dominance} bounds its
metric below by a positive constant times $g$.  The metric is smooth and
has all-order differentiated decay of order $2-n$ in normalized end
coordinates.  Thus the metric and tensor satisfy the stronger falloffs in
Theorem~\ref{thm:neutral-input}, with any exponent strictly between
$(n-2)/2$ and $n-2$.  The neutral DEC follows from
Lemma~\ref{flux:trace-tensor}; its matter densities are integrable because
the tensor is compactly supported and the scalar curvature is integrable
at infinity.  The ADM momentum is zero.  The ADM energy exists and is
given by \eqref{flux:mass-decrease}.  Finally the new boundary is strictly
future outer trapped by \eqref{flux:strict-new-boundary}.

Every full cut of $\mathcal D_N$, regarded as a subset of the original
exterior, is also a full cut there.  This includes any portions coincident
with $\partial\mathcal D_N$, and counts every component.  Therefore its
full-cut infimum $a_N$ satisfies
\begin{equation}
 a_N\geq e^{ky_0}a_*>0.
 \label{flux:area-retention}
\end{equation}
Indeed the metric comparison acts on all $k$ tangent directions, and the
class of available cuts has only been restricted.  No minimizing boundary
is assumed to be smooth or to exist.

It remains to verify the timelike energy hypothesis of the neutral input
for these auxiliary data.  Write $e_N=\mathcal E(\gamma)$.  If $e_N>0$,
their zero momentum makes that hypothesis immediate.  If $e_N\leq0$,
apply the conformal end increase of Lemma~\ref{prep:end-increase}, taking
its positive function $\phi$ constant on a core containing the boundary
and the new tensor's support.  For $U=1+\lambda\phi$, transform the data by
\begin{equation}
 \gamma_\lambda=U^{4/(n-2)}\gamma,
 \qquad K_\lambda=U^{2/(n-2)}K_{\mathrm{new}}.
 \label{flux:neutral-end-increase}
\end{equation}
The neutral DEC and trapping signs persist, the stronger asymptotic
derivatives persist, momentum remains zero, and full-cut areas do not
decrease.  Its energy is $e_N+2\lambda$.  For every arbitrarily small
$e>0$ choose $\lambda=(e-e_N)/2\geq0$.  These data satisfy all hypotheses
of Theorem~\ref{thm:neutral-input}, which would give
\begin{equation}
 e\geq\frac12(a_N/\omega)^{(k-1)/k}
 \geq\frac12e^{(k-1)y_0}X_*>0.
 \label{flux:nonpositive-contradiction}
\end{equation}
Letting $e\downarrow0$ is a contradiction.  Thus $e_N>0$, and the same
neutral input applied directly gives
\begin{equation}
 e^{(k-1)\varepsilon}(\mathcal E_*+o_N(1))
 -\frac{H_0(0)}k|Q|
 \geq\frac12 e^{(k-1)y_0}X_*
 =\frac{1-s^2}{2}X_*.
 \label{flux:prelimit-bound}
\end{equation}

The parameter order is as follows.  For the fixed prepared data, fix
$s\in(0,1)$, a profile tolerance, and $\varepsilon>0$ first.  Complete the
two scalar solves and the regular cut for each sufficiently large $N$;
the tensor plateau is chosen last for that $N$.  Now let $N\to\infty$ in
\eqref{flux:prelimit-bound}, then let $\varepsilon\downarrow0$, and then
let the profile tolerance decrease to zero.  Equations
\eqref{flux:profile-value} and \eqref{flux:floor} give
\eqref{flux:prepared-inequality}.  Only fixed-parameter estimates were
required in the second scalar solve.  The electromagnetic input used in
this proof was precisely \eqref{flux:graph-form-norm} and its lower bound
by $I_g(w)$, which also holds for the indicated compact magnetic form.
\end{proof}

\subsection{Return to the original data and optimize}

\begin{corollary}[Strict compact-magnetic numerical extension]
\label{flux:magnetic-extension}
Let the original exterior satisfy the numerical geometric and asymptotic
hypotheses, and let $D_m$ be an arbitrary smooth compactly supported closed
two-form.  With $\alpha=i_E\vol_g$ closed, use the electromagnetic convention
\begin{equation}
 I_g(v)=c_n^2\bigl(|E|_g^2+|D_m|_g^2+2D_m(v,E)\bigr),
 \qquad
 \mu_m+J_m(v)=\mu+J(v)-I_g(v).
 \label{flux:magnetic-convention}
\end{equation}
Assume also that $\mu_m,|J_m|_g$ are integrable and that the ADM and
charge flux limits exist with the stated normalizations.  Assume the ADM
vector is future timelike, the fixed componentwise trapping
inequalities are strict, and for some $c>0$ and
$0<\beta<\min(1,q)$,
\begin{equation}
 \mu_m-|J_m|_g\geq c r^{-n-\beta}.
 \label{flux:strict-magnetic-margin}
\end{equation}
Then, with the original full-cut infimum $a$ and the original normalized
electric charge $Q$,
\begin{equation}
 m\geq s|Q|+\frac{1-s^2}{2}X_A\quad(0<s<1).
 \label{flux:magnetic-s-family}
\end{equation}
Equivalently, $m\geq|Q|$, and when $X_A>|Q|$ one also has
$m\geq(X_A+Q^2/X_A)/2$.  No magnetic equality classification or
non-timelike extension is asserted here.
\end{corollary}

\begin{proof}
Apply the rest preparation and strict preparation in
Proposition~\ref{prep:rest-preparation} and
Lemma~\ref{prep:strictification}, retaining the two closed forms.  The
magnetic support stays in the compact original region.  There the strictly
positive margin absorbs the arbitrarily small rescaling errors.  The
gluing and bending annuli can be taken outside that support, so on those
annuli and on the end the electric estimates are unchanged.  This yields
prepared data with charge $Q$, energies $\mathcal E_{*,j}\to m$, and
full-cut infima $a_{*,j}\geq(1-o_j(1))a$.  For each fixed $s$, apply
Proposition~\ref{flux:prepared-bound} and then pass to this limit.  The
right side is increasing in the area variable, so the lower comparison of
infima is sufficient.  This proves \eqref{flux:magnetic-s-family}.  The
elementary optimization is given in
Proposition~\ref{prep:prepared-reduction}.  The argument does not require
the component electric fluxes to have one sign.
\end{proof}

\begin{proof}[Completion of Theorem~\ref{thm:numerical}]
Proposition~\ref{flux:prepared-bound} establishes
\eqref{prep:prepared-bound-hypothesis} for every prepared data set
with positive energy and zero momentum, and every fixed $s\in(0,1)$.
The Reduction of the Numerical Theorem
(Proposition~\ref{prep:prepared-reduction}) therefore proves
Theorem~\ref{thm:numerical}, including positivity of the full enclosing
area, strict future timelikeness of the original ADM vector, and the
stated branch distinction. The endpoint limits in $s$ are taken after
the fixed-$s$ inequality is established, so no profile estimate uniform
at $s=0$ or $s=1$ is required.
\end{proof}
